% --- Archon physics formalization source begin ---
% archon:physics
% archon:covers PhyXMiniProblems/problem_phyx_mini_0358.lean
% archon:source-report reports/phyx_mini/problem_phyx_mini_0358.source.json

\chapter{Physics problem phyx\_mini\_0358}
\label{ch:PhyXMiniProblems_problem_phyx_mini_0358}

\paragraph{Problem source.}
\#\# Physical scenario

An ideal diatomic gas has a molar internal energy equal to \textbackslash{}( E = \textbackslash{}frac\{5\}\{2\}RT \textbackslash{}) which depends only on its absolute temperature \textbackslash{}( T \textbackslash{}). A mole of this gas is taken quasi\textbackslash{}-statically first from state \textbackslash{}( A \textbackslash{}) to state \textbackslash{}( B \textbackslash{}), and then from state \textbackslash{}( B \textbackslash{}) to state \textbackslash{}( C \textbackslash{}) along the straight line paths shown in the diagram of pressure \textbackslash{}( p \textbackslash{}) versus volume \textbackslash{}( V \textbackslash{}).

\#\# Question

What is the work done by the gas in the process \textbackslash{}( A \textbackslash{}rightarrow B \textbackslash{}rightarrow C \textbackslash{})?

\#\# Answer choices

- A: A: 1300J

- B: B: 1895J

- C: C: 1550J

- D: D: 1600J

\#\# Figure caption (auxiliary; use the image as primary evidence)

The image is a graph showing a pressure-volume (p-V) diagram with three key points labeled A, B, and C. 

- **Axes**: 
  - The y-axis is labeled as \textbackslash{}( p \textbackslash{}) with units \textbackslash{}((10\textasciicircum{}\{6\} \textbackslash{}, \textbackslash{}text\{dynes\} \textbackslash{}, \textbackslash{}text\{cm\}\textasciicircum{}\{-2\})\textbackslash{}).
  - The x-axis is labeled as \textbackslash{}( V \textbackslash{}) with units \textbackslash{}((10\textasciicircum{}\{3\} \textbackslash{}, \textbackslash{}text\{cm\}\textasciicircum{}\{3\})\textbackslash{}).

- **Grid and Scale**:
  - Both axes have grids with equal spacing, allowing for accurate interpretation of values.
  
- **Points and Lines**:
  - Point A is around (1, 4) on the graph.
  - Point B is approximately at (2, 7).
  - Point C is roughly at (3, 5).
  - The points are connected by lines with arrows indicating the direction from A to B and B to C, suggesting a process path or sequence.

- **Text**: 
  - This includes the labels for the axes and the annotation of points A, B, and C.

The graph likely represents a thermodynamic process or cycle involving changes in pressure and volume.

\#\# Dataset metadata

Laws of Thermodynamics; Physical Model Grounding Reasoning, Multi-Formula Reasoning

\paragraph{Recorded answer/context.}
A: A: 1300J

\paragraph{Figure/image path.}
/root/proposal\_for\_physic/science-mango/phyx\_mini\_run/phyx\_data/test\_image/358.png

\paragraph{Formalization target.}
create a compiling Lean file with sorry bodies at `PhyXMiniProblems/problem\_phyx\_mini\_0358.lean`. The Lean declarations must preserve the physical quantities, dimensions or dimensional roles, figure labels, governing-law hypotheses, and final relation expressed by this problem.
Use Mathlib/Physlib names found through LeanExplore where available. If a physics API is missing, introduce faithful local abstractions rather than scalar placeholder aliases.

\begin{theorem}[Physics formalization target]
\label{thm:physics:phyx_mini_0358:target}
\lean{PhyXMiniProblems.ProblemPhyXMini0358.workDoneByGasAlongABC_eq_1300_joules}
\uses{def:physics:phyx-mini-0358:phyxminiproblems-problemphyxmini0358-diatomicgaspvprocess, def:physics:phyx-mini-0358:phyxminiproblems-problemphyxmini0358-matchesproblemandprimaryfigure, def:physics:phyx-mini-0358:phyxminiproblems-problemphyxmini0358-hasphysicalgasparameters, def:physics:phyx-mini-0358:phyxminiproblems-problemphyxmini0358-satisfiesdiatomicidealgascaloriclaw, def:physics:phyx-mini-0358:phyxminiproblems-problemphyxmini0358-satisfiesquasistaticstraightpathworklaw, def:physics:phyx-mini-0358:phyxminiproblems-problemphyxmini0358-totalworkbygasinjoules}
The assigned autoformalize agent should translate this physics problem into Lean declarations in the covered file, with theorem and lemma proof bodies written as `by sorry`.
\end{theorem}
\begin{proof}
This is an autoformalization task, not a proof task. Produce statements that can later be proved by the physics prover without weakening the physical model.
\end{proof}

% --- Archon declaration topology begin ---
\section{Declaration topology}
\begin{definition}[Gas Volume]
  \label{def:physics:phyx-mini-0358:phyxminiproblems-problemphyxmini0358-gasvolume}
  \lean{PhyXMiniProblems.ProblemPhyXMini0358.GasVolume}
  A signed physical gas volume, with dimension L³
\end{definition}
\begin{proof}
  This is the declared model definition of Gas Volume.
\end{proof}
\begin{definition}[Absolute Temperature]
  \label{def:physics:phyx-mini-0358:phyxminiproblems-problemphyxmini0358-absolutetemperature}
  \lean{PhyXMiniProblems.ProblemPhyXMini0358.AbsoluteTemperature}
  An absolute-temperature quantity, with Physlib's temperature dimension
\end{definition}
\begin{proof}
  This is the declared model definition of Absolute Temperature.
\end{proof}
\begin{definition}[Molar Gas Constant]
  \label{def:physics:phyx-mini-0358:phyxminiproblems-problemphyxmini0358-molargasconstant}
  \lean{PhyXMiniProblems.ProblemPhyXMini0358.MolarGasConstant}
  The molar gas constant, represented with energy-per-temperature dimension. Physlib's current Dimension has no amount-of-substance component, so the inverse-mole role is recorded by the name and paired with a scalar mole readout below
\end{definition}
\begin{proof}
  This is the declared model definition of Molar Gas Constant.
\end{proof}
\begin{definition}[Molar Internal Energy]
  \label{def:physics:phyx-mini-0358:phyxminiproblems-problemphyxmini0358-molarinternalenergy}
  \lean{PhyXMiniProblems.ProblemPhyXMini0358.MolarInternalEnergy}
  Molar internal energy. Its omitted inverse-mole component is handled in the same way as for MolarGasConstant; the underlying energy dimension is the grounded Physlib type DimEnergy
\end{definition}
\begin{proof}
  This is the declared model definition of Molar Internal Energy.
\end{proof}
\begin{definition}[volume In Cubic Meters]
  \label{def:physics:phyx-mini-0358:phyxminiproblems-problemphyxmini0358-volumeincubicmeters}
  \lean{PhyXMiniProblems.ProblemPhyXMini0358.volumeInCubicMeters}
  \uses{def:physics:phyx-mini-0358:phyxminiproblems-problemphyxmini0358-gasvolume}
  Read a physical volume in SI cubic metres
\end{definition}
\begin{proof}
  This is the declared model definition of volume In Cubic Meters.
\end{proof}
\begin{definition}[volume In Diagram Units]
  \label{def:physics:phyx-mini-0358:phyxminiproblems-problemphyxmini0358-volumeindiagramunits}
  \lean{PhyXMiniProblems.ProblemPhyXMini0358.volumeInDiagramUnits}
  \uses{def:physics:phyx-mini-0358:phyxminiproblems-problemphyxmini0358-gasvolume, def:physics:phyx-mini-0358:phyxminiproblems-problemphyxmini0358-volumeincubicmeters}
  Read a volume in the horizontal-axis unit 10³ cm³, equivalently litres
\end{definition}
\begin{proof}
  This is the declared model definition of volume In Diagram Units.
\end{proof}
\begin{definition}[pressure In Pascals]
  \label{def:physics:phyx-mini-0358:phyxminiproblems-problemphyxmini0358-pressureinpascals}
  \lean{PhyXMiniProblems.ProblemPhyXMini0358.pressureInPascals}
  Read a pressure in SI pascals
\end{definition}
\begin{proof}
  This is the declared model definition of pressure In Pascals.
\end{proof}
\begin{definition}[pressure In Diagram Units]
  \label{def:physics:phyx-mini-0358:phyxminiproblems-problemphyxmini0358-pressureindiagramunits}
  \lean{PhyXMiniProblems.ProblemPhyXMini0358.pressureInDiagramUnits}
  Read a pressure in the vertical-axis unit 10⁶ dyn cm⁻², which is exactly one bar or 10⁵ Pa
\end{definition}
\begin{proof}
  This is the declared model definition of pressure In Diagram Units.
\end{proof}
\begin{definition}[temperature In Kelvins]
  \label{def:physics:phyx-mini-0358:phyxminiproblems-problemphyxmini0358-temperatureinkelvins}
  \lean{PhyXMiniProblems.ProblemPhyXMini0358.temperatureInKelvins}
  \uses{def:physics:phyx-mini-0358:phyxminiproblems-problemphyxmini0358-absolutetemperature}
  Kelvin readout of a dimensionful absolute temperature
\end{definition}
\begin{proof}
  This is the declared model definition of temperature In Kelvins.
\end{proof}
\begin{definition}[molar Gas Constant In Joules Per Mole Kelvin]
  \label{def:physics:phyx-mini-0358:phyxminiproblems-problemphyxmini0358-molargasconstantinjoulespermolekelvin}
  \lean{PhyXMiniProblems.ProblemPhyXMini0358.molarGasConstantInJoulesPerMoleKelvin}
  \uses{def:physics:phyx-mini-0358:phyxminiproblems-problemphyxmini0358-molargasconstant}
  SI readout of the molar gas constant in joules per mole-kelvin
\end{definition}
\begin{proof}
  This is the declared model definition of molar Gas Constant In Joules Per Mole Kelvin.
\end{proof}
\begin{definition}[molar Internal Energy In Joules Per Mole]
  \label{def:physics:phyx-mini-0358:phyxminiproblems-problemphyxmini0358-molarinternalenergyinjoulespermole}
  \lean{PhyXMiniProblems.ProblemPhyXMini0358.molarInternalEnergyInJoulesPerMole}
  \uses{def:physics:phyx-mini-0358:phyxminiproblems-problemphyxmini0358-molarinternalenergy}
  Joule-per-mole readout of a molar internal energy
\end{definition}
\begin{proof}
  This is the declared model definition of molar Internal Energy In Joules Per Mole.
\end{proof}
\begin{definition}[work In Joules]
  \label{def:physics:phyx-mini-0358:phyxminiproblems-problemphyxmini0358-workinjoules}
  \lean{PhyXMiniProblems.ProblemPhyXMini0358.workInJoules}
  Joule readout of a signed physical work quantity
\end{definition}
\begin{proof}
  This is the declared model definition of work In Joules.
\end{proof}
\begin{definition}[pressure Axis Scale Dynes Per Square Centimeter]
  \label{def:physics:phyx-mini-0358:phyxminiproblems-problemphyxmini0358-pressureaxisscaledynespersquarecentimeter}
  \lean{PhyXMiniProblems.ProblemPhyXMini0358.pressureAxisScaleDynesPerSquareCentimeter}
  Number of dynes per square centimetre represented by one vertical unit
\end{definition}
\begin{proof}
  This is the declared model definition of pressure Axis Scale Dynes Per Square Centimeter.
\end{proof}
\begin{definition}[volume Axis Scale Cubic Centimeters]
  \label{def:physics:phyx-mini-0358:phyxminiproblems-problemphyxmini0358-volumeaxisscalecubiccentimeters}
  \lean{PhyXMiniProblems.ProblemPhyXMini0358.volumeAxisScaleCubicCentimeters}
  Number of cubic centimetres represented by one horizontal unit
\end{definition}
\begin{proof}
  This is the declared model definition of volume Axis Scale Cubic Centimeters.
\end{proof}
\begin{definition}[pressure Diagram Unit In Pascals]
  \label{def:physics:phyx-mini-0358:phyxminiproblems-problemphyxmini0358-pressurediagramunitinpascals}
  \lean{PhyXMiniProblems.ProblemPhyXMini0358.pressureDiagramUnitInPascals}
  SI pressure represented by one vertical plot unit
\end{definition}
\begin{proof}
  This is the declared model definition of pressure Diagram Unit In Pascals.
\end{proof}
\begin{definition}[volume Diagram Unit In Cubic Meters]
  \label{def:physics:phyx-mini-0358:phyxminiproblems-problemphyxmini0358-volumediagramunitincubicmeters}
  \lean{PhyXMiniProblems.ProblemPhyXMini0358.volumeDiagramUnitInCubicMeters}
  SI volume represented by one horizontal plot unit
\end{definition}
\begin{proof}
  This is the declared model definition of volume Diagram Unit In Cubic Meters.
\end{proof}
\begin{definition}[joules Per Diagram Area Unit]
  \label{def:physics:phyx-mini-0358:phyxminiproblems-problemphyxmini0358-joulesperdiagramareaunit}
  \lean{PhyXMiniProblems.ProblemPhyXMini0358.joulesPerDiagramAreaUnit}
  \uses{def:physics:phyx-mini-0358:phyxminiproblems-problemphyxmini0358-pressurediagramunitinpascals, def:physics:phyx-mini-0358:phyxminiproblems-problemphyxmini0358-volumediagramunitincubicmeters}
  Energy represented by one unit of area in the supplied pV diagram
\end{definition}
\begin{proof}
  This is the declared model definition of joules Per Diagram Area Unit.
\end{proof}
\begin{definition}[Gas Model]
  \label{def:physics:phyx-mini-0358:phyxminiproblems-problemphyxmini0358-gasmodel}
  \lean{PhyXMiniProblems.ProblemPhyXMini0358.GasModel}
  Material model specified in the prose
\end{definition}
\begin{proof}
  This is the declared model definition of Gas Model.
\end{proof}
\begin{definition}[PVDiagram State]
  \label{def:physics:phyx-mini-0358:phyxminiproblems-problemphyxmini0358-pvdiagramstate}
  \lean{PhyXMiniProblems.ProblemPhyXMini0358.PVDiagramState}
  The three state labels printed in the diagram
\end{definition}
\begin{proof}
  This is the declared model definition of PVDiagram State.
\end{proof}
\begin{definition}[Process Leg]
  \label{def:physics:phyx-mini-0358:phyxminiproblems-problemphyxmini0358-processleg}
  \lean{PhyXMiniProblems.ProblemPhyXMini0358.ProcessLeg}
  The two directed legs of the route A → B → C
\end{definition}
\begin{proof}
  This is the declared model definition of Process Leg.
\end{proof}
\begin{definition}[Process Regime]
  \label{def:physics:phyx-mini-0358:phyxminiproblems-problemphyxmini0358-processregime}
  \lean{PhyXMiniProblems.ProblemPhyXMini0358.ProcessRegime}
  Thermodynamic regime stated for both legs
\end{definition}
\begin{proof}
  This is the declared model definition of Process Regime.
\end{proof}
\begin{definition}[Path Shape]
  \label{def:physics:phyx-mini-0358:phyxminiproblems-problemphyxmini0358-pathshape}
  \lean{PhyXMiniProblems.ProblemPhyXMini0358.PathShape}
  Geometric shape of a leg in the pressure-volume plane
\end{definition}
\begin{proof}
  This is the declared model definition of Path Shape.
\end{proof}
\begin{definition}[Axis Label]
  \label{def:physics:phyx-mini-0358:phyxminiproblems-problemphyxmini0358-axislabel}
  \lean{PhyXMiniProblems.ProblemPhyXMini0358.AxisLabel}
  Literal variable labels appearing on the axes
\end{definition}
\begin{proof}
  This is the declared model definition of Axis Label.
\end{proof}
\begin{definition}[Axis Scale Label]
  \label{def:physics:phyx-mini-0358:phyxminiproblems-problemphyxmini0358-axisscalelabel}
  \lean{PhyXMiniProblems.ProblemPhyXMini0358.AxisScaleLabel}
  Literal scale labels printed beside the two axes
\end{definition}
\begin{proof}
  This is the declared model definition of Axis Scale Label.
\end{proof}
\begin{definition}[Diatomic Gas PVProcess]
  \label{def:physics:phyx-mini-0358:phyxminiproblems-problemphyxmini0358-diatomicgaspvprocess}
  \lean{PhyXMiniProblems.ProblemPhyXMini0358.DiatomicGasPVProcess}
  \uses{def:physics:phyx-mini-0358:phyxminiproblems-problemphyxmini0358-gasvolume, def:physics:phyx-mini-0358:phyxminiproblems-problemphyxmini0358-absolutetemperature, def:physics:phyx-mini-0358:phyxminiproblems-problemphyxmini0358-molargasconstant, def:physics:phyx-mini-0358:phyxminiproblems-problemphyxmini0358-molarinternalenergy, def:physics:phyx-mini-0358:phyxminiproblems-problemphyxmini0358-gasmodel, def:physics:phyx-mini-0358:phyxminiproblems-problemphyxmini0358-pvdiagramstate, def:physics:phyx-mini-0358:phyxminiproblems-problemphyxmini0358-processleg, def:physics:phyx-mini-0358:phyxminiproblems-problemphyxmini0358-processregime, def:physics:phyx-mini-0358:phyxminiproblems-problemphyxmini0358-pathshape, def:physics:phyx-mini-0358:phyxminiproblems-problemphyxmini0358-axislabel, def:physics:phyx-mini-0358:phyxminiproblems-problemphyxmini0358-axisscalelabel}
  The one-mole gas sample and its piecewise-straight process. The leg works are independent dimensionful fields. In particular, neither they nor their sum is assigned any requested numerical value here
\end{definition}
\begin{proof}
  This is the declared model definition of Diatomic Gas PVProcess.
\end{proof}
\begin{definition}[Matches Problem And Primary Figure]
  \label{def:physics:phyx-mini-0358:phyxminiproblems-problemphyxmini0358-matchesproblemandprimaryfigure}
  \lean{PhyXMiniProblems.ProblemPhyXMini0358.MatchesProblemAndPrimaryFigure}
  \uses{def:physics:phyx-mini-0358:phyxminiproblems-problemphyxmini0358-volumeindiagramunits, def:physics:phyx-mini-0358:phyxminiproblems-problemphyxmini0358-pressureindiagramunits, def:physics:phyx-mini-0358:phyxminiproblems-problemphyxmini0358-diatomicgaspvprocess}
  Problem-text data and exact transcription of the primary raster. The auxiliary caption's approximate (1,4), (2,7), (3,5) transcription is not used because the problem designates the image as primary evidence
\end{definition}
\begin{proof}
  This is the declared model definition of Matches Problem And Primary Figure.
\end{proof}
\begin{definition}[Has Physical Gas Parameters]
  \label{def:physics:phyx-mini-0358:phyxminiproblems-problemphyxmini0358-hasphysicalgasparameters}
  \lean{PhyXMiniProblems.ProblemPhyXMini0358.HasPhysicalGasParameters}
  \uses{def:physics:phyx-mini-0358:phyxminiproblems-problemphyxmini0358-volumeincubicmeters, def:physics:phyx-mini-0358:phyxminiproblems-problemphyxmini0358-pressureinpascals, def:physics:phyx-mini-0358:phyxminiproblems-problemphyxmini0358-temperatureinkelvins, def:physics:phyx-mini-0358:phyxminiproblems-problemphyxmini0358-molargasconstantinjoulespermolekelvin, def:physics:phyx-mini-0358:phyxminiproblems-problemphyxmini0358-diatomicgaspvprocess}
  Positivity conditions selecting physically meaningful gas states
\end{definition}
\begin{proof}
  This is the declared model definition of Has Physical Gas Parameters.
\end{proof}
\begin{definition}[Satisfies Diatomic Ideal Gas Caloric Law]
  \label{def:physics:phyx-mini-0358:phyxminiproblems-problemphyxmini0358-satisfiesdiatomicidealgascaloriclaw}
  \lean{PhyXMiniProblems.ProblemPhyXMini0358.SatisfiesDiatomicIdealGasCaloricLaw}
  \uses{def:physics:phyx-mini-0358:phyxminiproblems-problemphyxmini0358-temperatureinkelvins, def:physics:phyx-mini-0358:phyxminiproblems-problemphyxmini0358-molargasconstantinjoulespermolekelvin, def:physics:phyx-mini-0358:phyxminiproblems-problemphyxmini0358-molarinternalenergyinjoulespermole, def:physics:phyx-mini-0358:phyxminiproblems-problemphyxmini0358-diatomicgaspvprocess}
  The caloric equation of state stated in the problem: the molar internal energy is a function only of absolute temperature and obeys Uₘ(T) = (5/2) R T. This law makes no assertion about work
\end{definition}
\begin{proof}
  This is the declared model definition of Satisfies Diatomic Ideal Gas Caloric Law.
\end{proof}
\begin{definition}[Satisfies Quasi Static Straight Path Work Law]
  \label{def:physics:phyx-mini-0358:phyxminiproblems-problemphyxmini0358-satisfiesquasistaticstraightpathworklaw}
  \lean{PhyXMiniProblems.ProblemPhyXMini0358.SatisfiesQuasiStaticStraightPathWorkLaw}
  \uses{def:physics:phyx-mini-0358:phyxminiproblems-problemphyxmini0358-volumeindiagramunits, def:physics:phyx-mini-0358:phyxminiproblems-problemphyxmini0358-pressureindiagramunits, def:physics:phyx-mini-0358:phyxminiproblems-problemphyxmini0358-workinjoules, def:physics:phyx-mini-0358:phyxminiproblems-problemphyxmini0358-joulesperdiagramareaunit, def:physics:phyx-mini-0358:phyxminiproblems-problemphyxmini0358-diatomicgaspvprocess}
  Boundary-work law for a quasi-static straight segment in a pV diagram. Linear pressure variation makes the integral ∫ p dV equal average pressure times volume change. The final factor converts diagram-area units to joules. This general law contains none of the requested segment or total work values
\end{definition}
\begin{proof}
  This is the declared model definition of Satisfies Quasi Static Straight Path Work Law.
\end{proof}
\begin{definition}[total Work By Gas In Joules]
  \label{def:physics:phyx-mini-0358:phyxminiproblems-problemphyxmini0358-totalworkbygasinjoules}
  \lean{PhyXMiniProblems.ProblemPhyXMini0358.totalWorkByGasInJoules}
  \uses{def:physics:phyx-mini-0358:phyxminiproblems-problemphyxmini0358-workinjoules, def:physics:phyx-mini-0358:phyxminiproblems-problemphyxmini0358-diatomicgaspvprocess}
  Total work by the gas along the concatenated route A → B → C
\end{definition}
\begin{proof}
  This is the declared model definition of total Work By Gas In Joules.
\end{proof}
\begin{definition}[Answer Choice]
  \label{def:physics:phyx-mini-0358:phyxminiproblems-problemphyxmini0358-answerchoice}
  \lean{PhyXMiniProblems.ProblemPhyXMini0358.AnswerChoice}
  Labels of the four displayed answer choices
\end{definition}
\begin{proof}
  This is the declared model definition of Answer Choice.
\end{proof}
\begin{definition}[displayed Work In Joules]
  \label{def:physics:phyx-mini-0358:phyxminiproblems-problemphyxmini0358-displayedworkinjoules}
  \lean{PhyXMiniProblems.ProblemPhyXMini0358.displayedWorkInJoules}
  \uses{def:physics:phyx-mini-0358:phyxminiproblems-problemphyxmini0358-answerchoice}
  Work readout in joules printed beside each answer label
\end{definition}
\begin{proof}
  This is the declared model definition of displayed Work In Joules.
\end{proof}
\begin{definition}[recorded Answer Choice]
  \label{def:physics:phyx-mini-0358:phyxminiproblems-problemphyxmini0358-recordedanswerchoice}
  \lean{PhyXMiniProblems.ProblemPhyXMini0358.recordedAnswerChoice}
  \uses{def:physics:phyx-mini-0358:phyxminiproblems-problemphyxmini0358-answerchoice}
  The answer label recorded in the supplied dataset metadata
\end{definition}
\begin{proof}
  This is the declared model definition of recorded Answer Choice.
\end{proof}
\begin{lemma}[work Along Straight Process Legs]
  \label{lem:physics:phyx-mini-0358:phyxminiproblems-problemphyxmini0358-workalongstraightprocesslegs}
  \lean{PhyXMiniProblems.ProblemPhyXMini0358.workAlongStraightProcessLegs}
  \uses{def:physics:phyx-mini-0358:phyxminiproblems-problemphyxmini0358-workinjoules, def:physics:phyx-mini-0358:phyxminiproblems-problemphyxmini0358-diatomicgaspvprocess, def:physics:phyx-mini-0358:phyxminiproblems-problemphyxmini0358-matchesproblemandprimaryfigure, def:physics:phyx-mini-0358:phyxminiproblems-problemphyxmini0358-satisfiesquasistaticstraightpathworklaw}
  The two straight legs contribute 700 J and 600 J, respectively. Both equalities are conclusions derived from the general boundary-work law and the primary-image readouts; neither occurs in a premise structure
\end{lemma}
\begin{proof}
  The conclusion follows from the governing relations and figure readouts named in the statement of work Along Straight Process Legs.
\end{proof}
% --- Archon declaration topology end ---
% --- Archon physics formalization source end ---
